\documentclass[11pt,a4paper]{article}
\usepackage[margin=1in]{geometry}
\usepackage{amsmath,amssymb,amsthm,mathtools,booktabs,microtype,xcolor}
\usepackage[numbers,sort&compress]{natbib}
\usepackage[colorlinks=true,linkcolor=blue!50!black,citecolor=blue!50!black,urlcolor=blue!50!black]{hyperref}
\newtheorem{theorem}{Theorem}[section]
\newtheorem{proposition}[theorem]{Proposition}
\newtheorem{corollary}[theorem]{Corollary}
\theoremstyle{remark}\newtheorem{remark}[theorem]{Remark}
\newcommand{\End}{\operatorname{End}}
\newcommand{\Tr}{\operatorname{Tr}}
\newcommand{\ad}{\operatorname{ad}}
\newcommand{\Hj}{\mathcal H_j}
\title{Half-Turn Covariance of a Tangential Fuzzy-Sphere Oscillator\\and Its Sector-Resolved Planar Limit}
\author{}
\date{29 October 2026}
\hypersetup{pdfauthor={},pdftitle={Half-Turn Covariance of a Tangential Fuzzy-Sphere Oscillator and Its Sector-Resolved Planar Limit}}
\begin{document}
\maketitle
\begin{abstract}
We study a positive oscillator operator on the complex matrix algebra of a spin-$j$ representation, using only the two components tangent to a chosen polar axis.
Its self-dual spectrum and parameter-inversion covariance follow from elementary left--right multiplication identities.
We give the exact curvature defect on each fixed polar occupation window, distinguish this statement from the full low-energy spectrum, and prove that the full self-dual heat trace converges to four copies of the planar oscillator heat trace.
Two cross-pole matrix units remain exact finite-energy eigenvectors for every oscillator parameter, so global spectral decoupling cannot be inferred from a selected-pole contraction.
We also derive the Weyl-symbol formula for right multiplication by parity: with the stated convention it is an involutive scaled symplectic Fourier transform.
A comparison with the earlier fuzzy harmonic regulator identifies the extra axial commutator term and prevents a broad novelty claim for fuzzy Grosse--Wulkenhaar actions.
The results concern a finite kinetic family and its explicit spectral limits; they do not assert interacting renormalization or operator-topology convergence.
\end{abstract}

\section{The finite question and its context}
Coordinate--momentum covariance is a defining feature of the noncommutative oscillator models associated with Langmann--Szabo duality and Grosse--Wulkenhaar renormalization~\citep{LangmannSzabo2002,GrosseWulkenhaar2005}.
Finite fuzzy-sphere regularizations and their planar limits are established tools~\citep{ChuMadoreSteinacker2001}.
In particular, Ydri constructed a fuzzy harmonic regulator and quartic action before the finite family considered here~\citep[Sec.~4.2]{Ydri2014}.
The useful question is therefore specific: what exact covariance and spectral information follow when the commutator and anticommutator terms are both restricted to two tangential components?

The answer has two parts.
The finite identities are exact and need no limiting argument.
The planar interpretation, however, must retain the multiplicity of the polar sectors and distinguish an ordinary Fourier transform from a scaled symplectic one.
We prove those statements directly and give portable exact checks of the finite conventions.
The contribution is the explicit analysis of this tangential specialization and its limits, rather than a new construction of fuzzy geometry or a first fuzzy oscillator action.

\section{Definition and positivity}
Let $j\in\{\tfrac12,1,\tfrac32,\ldots\}$, let $V_j$ be the spin-$j$ irreducible representation, and let $J_a$ be its Hermitian generators:
\begin{equation}
 [J_a,J_b]=i\epsilon_{abc}J_c,\qquad
 \sum_{a=1}^3J_a^2=j(j+1)I.
\end{equation}
The Hilbert space is $\Hj=\End(V_j)$ with inner product $\langle X,Y\rangle=\Tr(X^\dagger Y)$.
Define left and right multiplication by
\begin{equation}
 L_aX=J_aX,\qquad R_aX=XJ_a,\qquad
 A_a=L_a-R_a,\qquad S_a=L_a+R_a.
\end{equation}
All four operators are self-adjoint on $\Hj$, and left multiplication commutes with right multiplication.
Right multiplication has the opposite Lie-bracket convention, a fact that is harmless here because the calculations use explicit compositions.

For $\theta>0$ and $\Omega>0$, put
\begin{equation}\label{eq:H}
 H_j(\Omega)=\frac{1}{\theta j}\sum_{a=1}^2
 \big(A_a^2+\Omega^2S_a^2\big).
\end{equation}
It is positive because
\begin{equation}
 \langle X,H_j(\Omega)X\rangle
 =\frac{1}{\theta j}\sum_{a=1}^2
 \big(\|A_aX\|_{\rm HS}^2+\Omega^2\|S_aX\|_{\rm HS}^2\big)\ge0.
\end{equation}
We fix an axis once and for all; the index restriction $a=1,2$ is part of the model.
When $\theta$ is an area, $H_j$ has dimensions of inverse area and a heat parameter $t$ has dimensions of area.

The complex domain matters.
Although $H_j$ preserves Hermitian matrices, the right multiplication used below generally does not.
Its covariance is a unitary equivalence on the complex Hilbert space $\Hj$ and does not by itself define a symmetry of a real Hermitian field configuration space.

\section{Exact spectrum and covariance}
\begin{theorem}[Self-dual spectrum]\label{thm:spectrum}
Let $E_{m_Lm_R}=|m_L\rangle\langle m_R|$ with $m_L,m_R=-j,-j+1,\ldots,j$.
Then
\begin{equation}\label{eq:energy}
 H_j(1)E_{m_Lm_R}
 =\frac{2}{\theta j}\big(2j(j+1)-m_L^2-m_R^2\big)E_{m_Lm_R}.
\end{equation}
Equivalently, for $n_L=j-m_L$ and $n_R=j-m_R$,
\begin{equation}\label{eq:occupation}
 E_j(n_L,n_R)=\frac4\theta(n_L+n_R+1)
 -\frac2{\theta j}(n_L^2+n_R^2),\qquad 0\le n_L,n_R\le2j.
\end{equation}
\end{theorem}
\begin{proof}
For each $a$, $(L_a-R_a)^2+(L_a+R_a)^2=2(L_a^2+R_a^2)$.
The Casimir identity gives
\begin{equation}
 H_j(1)=\frac2{\theta j}
 \big(2j(j+1)I-L_3^2-R_3^2\big).
\end{equation}
The matrix units diagonalize $L_3$ and $R_3$.
Substituting $m=j-n$ yields~\eqref{eq:occupation}.
\end{proof}

\begin{theorem}[Finite parameter inversion]\label{thm:covariance}
Let $U_j=e^{i\pi J_3}$ and $\mathcal P_jX=XU_j$.
Then $\mathcal P_j$ is unitary,
\begin{equation}\label{eq:inversion}
 \mathcal P_jH_j(\Omega)\mathcal P_j^{-1}
 =\Omega^2H_j(\Omega^{-1}),\qquad
 \mathcal P_j^2=(-1)^{2j}I.
\end{equation}
\end{theorem}
\begin{proof}
The Hilbert--Schmidt norm is invariant under right multiplication by $U_j$.
Left multiplication is unchanged under conjugation by $\mathcal P_j$, whereas
\begin{equation}
 \mathcal P_jR_a\mathcal P_j^{-1}X
 =XU_j^{-1}J_aU_j=-R_aX,\qquad a=1,2.
\end{equation}
Thus $A_a$ and $S_a$ are exchanged, proving~\eqref{eq:inversion}.
Finally, $U_j^2=e^{2\pi iJ_3}=(-1)^{2j}I$.
\end{proof}

The spectrum at parameter $\Omega$ is therefore exactly $\Omega^2$ times the spectrum at $\Omega^{-1}$, including multiplicities.
A floating-point residual between these two spectra tests numerical accuracy, not a physical $1/j$ defect.
The curvature defect occurs in~\eqref{eq:occupation}, which compares a finite polar spectrum to a planar one.

\section{Polar windows and the full heat trace}
\begin{corollary}[Uniform fixed-window estimate]\label{cor:window}
Fix an integer $M\ge0$ and restrict $H_j(1)$ to the matrix units with $0\le n_L,n_R\le M$, for $2j\ge M$.
This subspace is invariant and its eigenvalues satisfy
\begin{equation}
 \left|E_j(n_L,n_R)-\frac4\theta(n_L+n_R+1)\right|
 \le\frac{4M^2}{\theta j}.
\end{equation}
\end{corollary}
\begin{proof}
The invariance follows from diagonalization and the bound follows from~\eqref{eq:occupation}.
\end{proof}

This estimate selects one corner of the occupation square.
The full spectrum contains other corners at the same energy scale.
For example, its ground energy is $4/\theta$ and its ground multiplicity is four for every $j>0$, whereas one planar oscillator has ground multiplicity one.

\begin{theorem}[Four-copy heat-trace limit]\label{thm:heat}
For every $t>0$,
\begin{equation}\label{eq:heatlimit}
 \lim_{j\to\infty}\Tr_{\Hj}e^{-tH_j(1)}
 =\frac{4e^{-4t/\theta}}{(1-e^{-4t/\theta})^2}.
\end{equation}
The limit holds through integer and half-integer spins.
\end{theorem}
\begin{proof}
For a magnetic number $m$, let $k=j-|m|$.
The possible $k$ are nonnegative integers up to $\lfloor j\rfloor$.
Their multiplicity $c_j(k)$ is two, except that $c_j(j)=1$ at the integer-spin midpoint.
Thus~\eqref{eq:energy} becomes
\begin{equation}
 E_j(k,l)=\frac4\theta+\frac4\theta(k+l)
 -\frac2{\theta j}(k^2+l^2).
\end{equation}
Since $0\le k,l\le j$, this energy is at least $4/\theta+2(k+l)/\theta$.
Extend $c_j(k)$ by zero outside its range.
The heat trace is the double sum of $c_j(k)c_j(l)e^{-tE_j(k,l)}$, dominated by the summable sequence
\begin{equation}
 4e^{-4t/\theta}e^{-2t(k+l)/\theta}.
\end{equation}
For every fixed $k,l$, $c_j(k)c_j(l)\to4$ and $E_j(k,l)\to4(k+l+1)/\theta$.
Dominated convergence and two geometric sums give~\eqref{eq:heatlimit}.
\end{proof}

The right-hand side is four times the heat trace of the planar self-dual matrix oscillator with occupation energies $4(n_L+n_R+1)/\theta$.
The factor four records the independent north/south choices on the ket and bra factors.
It is not removed by sending the matrix size to infinity.

\begin{proposition}[Persistent cross-pole states]\label{prop:crosspole}
For every $j>0$ and $\Omega>0$, the two matrix units $|j\rangle\langle-j|$ and $|-j\rangle\langle j|$ are eigenvectors with energy
\begin{equation}\label{eq:crosspole}
 \frac{2(1+\Omega^2)}{\theta}.
\end{equation}
\end{proposition}
\begin{proof}
Expanding~\eqref{eq:H} gives
\begin{equation}
 \theta jH_j(\Omega)
 =(1+\Omega^2)\sum_{a=1}^2(L_a^2+R_a^2)
 +2(\Omega^2-1)\sum_{a=1}^2L_aR_a.
\end{equation}
On either cross-pole unit, the first sum is $2j$ times the identity.
The second sum vanishes because
\begin{equation}
 \sum_{a=1}^2J_aXJ_a=\tfrac12(J_+XJ_-+J_-XJ_+)
\end{equation}
and each term contains an annihilated extremal ket or bra.
\end{proof}

Consequently, these states do not decouple by moving to infinite energy.
A single-pole continuum description may still be appropriate after selecting a sector or observable class, but that is additional structure.
Neither Theorem~\ref{thm:heat} nor Corollary~\ref{cor:window} defines embeddings between the varying matrix Hilbert spaces, so neither is an assertion of strong-, norm-, or resolvent convergence.

\section{Right parity and the symplectic Fourier convention}
The order of a transform is not enough to identify its action.
In particular, the ordinary Fourier transform and a Fourier transform composed with a symplectic rotation have different squares.
We give the relevant formula rather than importing an order-four analogy.

Let $A$ initially have a Schwartz integral kernel on $L^2(\mathbb R)$ and let $\Pi\psi(q)=\psi(-q)$.
Use the Weyl symbol
\begin{equation}\label{eq:weyl}
 a_A(q,p)=\int_{\mathbb R}e^{-ipy/\hbar}
 \langle q+y/2|A|q-y/2\rangle\,dy,\qquad\hbar>0.
\end{equation}
\begin{proposition}[Symbol of right parity]\label{prop:weyl}
In the convention~\eqref{eq:weyl},
\begin{equation}\label{eq:symplectic}
 a_{A\Pi}(q,p)=\frac1{\pi\hbar}\int_{\mathbb R^2}
 e^{2i(qp'-pq')/\hbar}a_A(q',p')\,dq'\,dp'.
\end{equation}
The transform in~\eqref{eq:symplectic} is unitary and involutive on $L^2(\mathbb R^2)$.
\end{proposition}
\begin{proof}
The inverse symbol formula is
\begin{equation}
 \langle u|A|v\rangle=\frac1{2\pi\hbar}
 \int e^{ip'(u-v)/\hbar}a_A((u+v)/2,p')\,dp'.
\end{equation}
Since $\langle u|A\Pi|v\rangle=\langle u|A|-v\rangle$, substituting $u=q+y/2$, $v=q-y/2$ gives center $y/2$ and separation $2q$.
The change of variable $q'=y/2$ proves~\eqref{eq:symplectic}.
Composing the skew kernel twice produces the delta distribution at the original $(q,p)$; equivalently it is a rescaled ordinary Fourier transform followed by the area-preserving symplectic rotation.
Plancherel gives unitarity with the displayed normalization.
The Schwartz identity extends by density to Hilbert--Schmidt operators and their $L^2$ symbols.
\end{proof}

In oscillator number states, $\Pi=(-1)^{\hat n}$.
On a fixed north-pole window,
\begin{equation}
 e^{i\pi J_3}|j-n\rangle=e^{i\pi j}(-1)^n|j-n\rangle.
\end{equation}
After removing the constant phase, the finite right half-turn therefore has exactly the right-parity action on these labels.
Proposition~\ref{prop:weyl} explains the appropriate planar Fourier interpretation of that action.
It also agrees with the involutive scaled Fourier duality in the original Langmann--Szabo convention~\citep[Prop.~1]{LangmannSzabo2002}.
This does not construct a preferred-axis-independent fuzzy-sphere Fourier transform, and it does not turn label agreement into an operator-topology contraction theorem.

\section{Comparison with the earlier fuzzy harmonic regulator}
The prior harmonic regulator must be compared at the operator level.
After choosing a common positive kinetic normalization, the squashed fuzzy-sphere operator of Ydri~\citep[Sec.~4.2]{Ydri2014} has dimensionless generator form
\begin{equation}
 Q_Y(\Omega)=\sum_{a=1}^3A_a^2+\Omega^2A_3^2
 +\Omega^2\sum_{a=1}^2S_a^2.
\end{equation}
The dimensionless numerator $Q_T=\theta jH_j$ of~\eqref{eq:H} instead satisfies
\begin{equation}\label{eq:prior}
 Q_Y(\Omega)-Q_T(\Omega)=(1+\Omega^2)A_3^2.
\end{equation}
Thus the two families differ by a positive axial commutator term.
They agree on the $A_3=0$ sector, but are not identical finite kinetic operators.
For example, the added term on $E_{m_Lm_R}$ is $(1+\Omega^2)(m_L-m_R)^2$.
The extra term changes the role of cross-pole states and cannot be discarded in a global spectral comparison.

Equation~\eqref{eq:prior} identifies a precise specialization worth studying while acknowledging the earlier fuzzy harmonic and quartic constructions.
No claim of a first fuzzy Grosse--Wulkenhaar action, interacting continuum renormalization, or publication priority is made.
The standard fuzzy-sphere and planar oscillator literature supplies the surrounding framework~\citep{ChuMadoreSteinacker2001,GrosseWulkenhaar2005}; the proofs here state the exact content of the tangential finite family.

\section{Reproducibility and interpretation}
The accompanying Python script constructs the spin generators independently using exact SymPy radicals.
It checks the complete self-dual superoperator, symbolic-$\Omega$ covariance, projective square, cross-pole eigenvectors, and~\eqref{eq:prior} for spins $1/2,1,3/2,2$.
It also records heat-trace sums approaching~\eqref{eq:heatlimit}.
Those decimals illustrate the limit; the proof is the summable domination above.
No Monte Carlo data, fitted spectrum or external project directory is needed to reproduce these checks.

The finite model is positive, explicitly solvable at its self-dual point, and exactly covariant under parameter inversion.
Its compact geometry leaves four polar sectors in the full thermal limit.
The useful continuum statement is therefore sector-resolved spectral convergence together with the all-sector heat-trace limit, rather than an unqualified identification of the full finite spectrum with one planar oscillator.

\section*{AI-assisted workflow and unresolved attribution}
The repository operator states:
\begin{quote}
I genuinely do not know how authorship, intellectual credit, or ownership should be attributed for these results.
\end{quote}
The operator reports that earlier work was led mainly by Fable 5, whereas results from July 2026 onward were predominantly produced by GPT models.
DeepSeek and multiple GPT models also participated in cross-review, checking, and manuscript preparation.
This is a reported chronology rather than an audited allocation of each individual contribution.
A later revision does not reattribute an earlier derivation.
Human author names and affiliations are intentionally blank, and the scientific claims have not undergone human expert review.
Historical research dates, the present edition date, private upload dates, and demonstrated public disclosure are distinct; no settled authorship, ownership, priority, or journal-acceptance claim is made.
\bibliographystyle{plainnat}
\bibliography{references}
\end{document}
